\section{Introducción}

Este es un ejemplo de documento con encabezados y varias fórmulas matemáticas.

\subsection{Ecuaciones básicas}

La ecuación de segundo grado puede escribirse como:

\[
ax^2 + bx + c = 0
\]

y su solución es:

\[
x = \frac{-b \pm \sqrt{b^2 - 4ac}}{2a}
\]

\subsection{Geometría}

El teorema de Pitágoras establece que:

\[
a^2 + b^2 = c^2
\]

El área de un círculo es:

\[
A = \pi r^2
\]

\section{Cálculo}

\subsection{Límites}

Un límite clásico es:

\[
\lim_{x \to \infty} \frac{1}{x} = 0
\]

Otro ejemplo:

\[
\lim_{x \to 0} \frac{\sin x}{x} = 1
\]

\subsection{Derivadas}

La derivada de una potencia es:

\[
\frac{d}{dx}(x^n) = nx^{n-1}
\]

Por ejemplo:

\[
\frac{d}{dx}(x^3) = 3x^2
\]

\subsection{Integrales}

Una integral sencilla es:

\[
\int_0^1 x^2 \, dx = \frac{1}{3}
\]

Y una integral indefinida:

\[
\int (3x^2 - 1)\,dx = x^3 - x + C
\]

\section{Series}

La serie geométrica infinita cumple:

\[
\sum_{n=0}^{\infty} r^n = \frac{1}{1-r}, \quad \text{si } |r| < 1
\]

La serie armónica es:

\[
\sum_{n=1}^{\infty} \frac{1}{n}
\]

y es divergente.

\section{Matrices}

Una matriz de ejemplo es:

\[
A =
\begin{pmatrix}
1 & 2 \\
3 & 4
\end{pmatrix}
\]

Su determinante vale:

\[
\det(A) = 1 \cdot 4 - 2 \cdot 3 = -2
\]
